
\newcommand{\ifTax}{\ifthenelse{\boolean{Taxes}}}

Consider a firm characterized by the following:
\begin{center}
\begin{tabular}{rcl}
   $\kap_{t}      $      & - & Firm's capital stock at the beginning of period $t$
\\ $\fFunc(\kap)     $ & - & The firm's total output depends only on $k$
\ifTax{\\ $\TaxCorp $ & - & Tax rate on corporate earnings
\\ $\TaxFree$  & - & $1-\TaxCorp$ = Portion of earnings untaxed
\\ $\rev_{t} = \fFunc(\kap_{t})\TaxFree   $ & - & After tax revenues
}{}
\\ $\inv_{t}      $ & - & Investment in period $t$
\\ $\jFunc(i,k)      $ & - & AdJustment costs associated with investment $\inv$ given capital $k$
\ifTax{\\ $\itc$    & - & Investment tax credit (ITC)
\\ $\kPriceAfterITC=\PostITC$    & - & $=1-\itc$ = Cost of 1 unit of (price-1) investment after ITC
\\ $\xpend_{t} = (\inv_{t}+\adj_{t})\kPriceAfterITC $  & - & after-tax eXpenditures (purchases plus adjustment costs) on investment
}{\\ $\xpend_{t} = \inv_{t}+\adj_{t}  $  & - & eXpenditures (purchases plus adjustment costs) on investment}
\\ $\Discount=1/\Rfree     $ & - & Discount factor for future profits (inverse of interest factor)
\end{tabular}
\end{center}

Suppose that the firm's goal is to pick the sequence $\iFunc_{t}$ that solves:
\begin{equation}\begin{gathered}\begin{aligned}
  \vFirm(\kap_{t}) & =  \max_{\{\iFunc\}_{t}^{\infty}}~\sum_{n=0}^{\infty} \Discount^{n} \left(\ifTax{\rev_{t+n}-\xpend_{t+n}}{f_{t+n}-\inv_{t+n}-\adj_{t+n}}\right)
\end{aligned}\end{gathered}\end{equation}
subject to the transition equation for capital,
\begin{equation}\begin{gathered}\begin{aligned}
  \kap_{t+1} & =  (\kap_{t}+\inv_{t})\DeprFac \label{eq:kAccum}
\end{aligned}\end{gathered}\end{equation}
where $\DeprFac = (1-\depr)$ is the amount of capital left after one
period of depreciation at rate $\depr$.\footnote{There are some small
  differences between the formulation of the model here and in
  \texttt{qModel}.  Here, investment costs are paid at the time of
  investment and the depreciation factor applies to $(\kap_{t}+\inv_{t})$
  rather than just $\kap_{t}$. These changes simplify the computational
  solution without changing any key results.}  $\vFirm_{t}$ is the
value of the profit-maximizing firm: If capital markets are efficient
this is the equity value that the firm would command if somebody
wanted to buy it.

\ifthenelse{\boolean{Exam}}{\begin{enumerate}}{}
\ifthenelse{\boolean{Exam}}
{\item Show that the Bellman equation for the firm can be derived as follows:
\begin{equation}\begin{gathered}\begin{aligned}
  \vFirm_{t}(\kap_{t}) & =  \max_{\{\inv_{t}\}}~ \ifTax{\rev_{t}-\xpend_{t}}{f_{t}-\inv_{t}-\adj_{t}}+\Discount \vFirm_{t+1}\left((\kap_{t}+\inv_{t})\DeprFac\right)
\end{aligned}\end{gathered}\end{equation}
}{The firm's Bellman equation can be written:}
\answer{
\begin{equation*}\begin{gathered}\begin{aligned}
  \vFirm_{t}(\kap_{t}) & =  \max_{\{\iFunc\}_{t}^{\infty}}~\sum_{n=0}^{\infty} \Discount^{n} \left(\ifTax{\rev_{t+n}-\xpend_{t+n}}{f_{t+n}-\inv_{t+n}-\adj_{t+n}}\right)
  \\ & =  \max_{\{\inv_{t}\}}~ \ifTax{\rev_{t}-\xpend_{t}}{f_{t}-\inv_{t}-\jFunc(\inv_{t},\kap_{t})}+\Discount \left[\max_{\{i\}_{t+1}^{\infty}}\sum_{n=0}^{\infty} \Discount^{n} \left(\ifTax{\rev_{t+1+n}-\xpend_{t+1+n}}{f_{t+1+n}-\inv_{t+1+n}-\adj_{t+1+n}}\right)\right]
\\ & =  \max_{\{\inv_{t}\}}~ \ifTax{\rev_{t}-\xpend_{t}}{f_{t}-\inv_{t}-\jFunc(\inv_{t},\kap_{t})}+\Discount \vFirm_{t+1}\left((\kap_{t}+\inv_{t})\DeprFac\right)
\end{aligned}\end{gathered}\end{equation*}
}{\AnswerOrBlankSpace{\clearpage\vfill\eject}{\vspace{0.0in}}}{}

\ifthenelse{\boolean{Exam}}{}{}
Define $\adj_{t}^{\inv}$ as the derivative of adjustment costs with respect to
the level of investment. \medskip

\ifthenelse{\boolean{Exam}}{\item Show that the first order condition for optimal investment implies
\begin{equation}\begin{gathered}\begin{aligned}
\ifTax{(1+j^{\inv}_{t})\kPriceAfterITC}{1+j^{\inv}_{t}} & =  \DeprFac\Discount \vFirm_{t+1}^{\kap}(\kap_{t+1}) \label{eq:iFOC}
\end{aligned}\end{gathered}\end{equation}
and provide a verbal interpretation of the condition.
}{
The first order condition for optimal investment implies:
}
\answer{
\begin{equation}\begin{gathered}\begin{aligned}
  0 & =  \ifTax{-\kPriceAfterITC(1+j^{\inv}_{t})}{-1 -j^{\inv}_{t} }+ \DeprFac\Discount \vFirm_{t+1}^{\kap}(\kap_{t+1})
\\ \ifTax{(1+j^{\inv}_{t})\kPriceAfterITC}{1+j^{\inv}_{t}} & =  \DeprFac\Discount \vFirm_{t+1}^{\kap}(\kap_{t+1}) \ifthenelse{\boolean{Exam}}{}{\label{eq:iFOC}}
\end{aligned}\end{gathered}\end{equation}

In words: The marginal \ifTax{after-tax}{} cost of an additional unit of
investment (the LHS) should be equal to the discounted marginal value
of the resulting extra capital (the RHS).


}{\AnswerOrBlankSpace{\clearpage\vfill\eject}{\vspace{0.0in}}}{}


\ifthenelse{\boolean{Exam}}
{\item Now use the Envelope theorem to derive the Euler equation for investment
\begin{equation}\begin{gathered}\begin{aligned}
 (1+\adj_{t}^{\inv})\ifTax{\kPriceAfterITC}{}  & = \DeprFac\Discount \left[ \ifTax{\TaxFree}{}\fFunc^{\kap}(\kap_{t+1})+(1+\adj_{t+1}^{\inv}-\adj_{t+1}^{\kap})\ifTax{\kPriceAfterITC}{}\right]
\label{eq:iEuler}
\end{aligned}\end{gathered}\end{equation}
}{
}

\answer{
The Envelope theorem says
\begin{equation}\begin{gathered}\begin{aligned}
  \vFirm_{t}^{\kap}(\kap_{t}) & =  \ifTax{\TaxFree}{}\fFunc^{\kap}(\kap_{t}) - \ifTax{\kPriceAfterITC}{}\adj_{t}^{\kap}+\Discount \vFirm_{t+1}^{\kap}(\kap_{t+1})\overbrace{\left(\frac{\partial \kap_{t+1}}{\partial \kap_{t}}\right)}^{\DeprFac} \nonumber
\\  & =  \ifTax{\TaxFree}{}\fFunc^{\kap}(\kap_{t}) - \ifTax{\kPriceAfterITC}{}\adj_{t}^{\kap}+\underbrace{\Discount \DeprFac \vFirm_{t+1}^{\kap}(\kap_{t+1})}_{=\ifTax{\kPriceAfterITC}{}(1+\adj_{t}^{\inv})~\text{from \eqref{eq:iFOC}}}
\label{eq:iEnvelope}
\end{aligned}\end{gathered}\end{equation}

So the corresponding $t+1$ equation can be substituted into \eqref{eq:iFOC}
to obtain
\begin{equation}\begin{gathered}\begin{aligned}
 (1+\adj_{t}^{\inv})\ifTax{\kPriceAfterITC}{}  & = \left( \ifTax{\TaxFree}{}\fFunc^{\kap}(\kap_{t+1})+(1+\adj_{t+1}^{\inv}-\adj_{t+1}^{\kap})\ifTax{\kPriceAfterITC}{}\right)\DeprFac\Discount
\ifthenelse{\boolean{Exam}}{}{\label{eq:iEuler}}
\end{aligned}\end{gathered}\end{equation}
which is the Euler equation for investment.

}{\AnswerOrBlankSpace{\clearpage\vfill\eject}{\vspace{0.0in}}}{}


Now suppose that a steady state exists in which the capital stock is
at its optimal level and is not adjusting, so costs of adjustment are
zero: $\adj_{t}=\adj_{t+1}=j^{\inv}_{t}=j^{\inv}_{t+1}=j^{\kap}_{t}=j^{\kap}_{t+1}=0$.


\ifthenelse{\boolean{Exam}}{
\item Use the investment Euler equation to show that the steady state level of the capital stock $\check{k}$
satisfies the equation below, and provide an intuitive interpretation of the equation.
\begin{equation}\begin{gathered}\begin{aligned}
 \ifTax{\kPriceAfterITC}{}\Rfree & =  \DeprFac (\ifTax{\kPriceAfterITC}{1}+\ifTax{\TaxFree}{}\fFunc^{\kap}(\check{k})) \label{eq:sskbar}
\end{aligned}\end{gathered}\end{equation}
}{}

\answer{
If $j^{\inv}_{t} = j^{\inv}_{t+1} = j^{\kap}_{t+1}$ then \eqref{eq:iEuler} reduces
to
\begin{equation}\begin{gathered}\begin{aligned}
 \ifTax{\kPriceAfterITC}{1} & = \overbrace{\Discount}^{=\Rfree^{-1}} \DeprFac\left[\ifTax{\TaxFree}{}\fFunc^{\kap}(\check{k})+\ifTax{\kPriceAfterITC}{1}\right]
\\ \ifTax{\kPriceAfterITC}{}\Rfree & =  \DeprFac (\ifTax{\kPriceAfterITC+}{}\ifTax{\TaxFree}{}\fFunc^{\kap}(\check{k}))\label{kSSimplicit}
\end{aligned}\end{gathered}\end{equation}
so that the capital stock is equal to the value that causes its \ifTax{after-tax}{}
marginal product to match the interest factor, after compensating for depreciation.
}{\AnswerOrBlankSpace{\clearpage\vfill\eject}{\vspace{0.0in}}}{}

Another way to analyze this problem is in terms of the marginal value of
capital, $\ek_{t} \equiv \vFirm_{t}^{\kap}(\kap_{t}).$

\ifthenelse{\boolean{Exam}}{
\item Show that in the vicinity of the steady state, assuming that adjustment
costs are approximately zero, the equation for $\ek$ will be
  \begin{equation}\begin{gathered}\begin{aligned}
\ek_{t} & =  \frac{\ifTax{\TaxFree}{}\fFunc^{\kap}(\kap_{t}) - \ifTax{\kPriceAfterITC}{}\adj_{t}^{\kap}+\Delta \ek_{t+1}}{(1-\Discount\DeprFac)}
  \end{aligned}\end{gathered}\end{equation}
  and use this equation along with the transition equation for capital
  to draw a phase diagram in $(\kap,\ek)$ space for this model.  Be
  sure to explain why the $\Delta \ek_{t+1}=0$ locus is downward
  sloping in the vicinity of the steady state.
}{}

%Assume for convenience that the depreciation rate is $\depr=0$.

\answer{
Rewrite \eqref{eq:iEnvelope} as
\begin{equation}\begin{gathered}\begin{aligned}
  \ek_{t} & =  \ifTax{\TaxFree}{}\fFunc^{\kap}(\kap_{t}) - \ifTax{\kPriceAfterITC}{}\adj_{t}^{\kap}+\Discount \DeprFac (\ek_{t}+\ek_{t+1}-\ek_{t})
\\ & =  \ifTax{\TaxFree}{}\fFunc^{\kap}(\kap_{t}) - \ifTax{\kPriceAfterITC}{}\adj_{t}^{\kap}+\Discount \DeprFac (\ek_{t}+\Delta \ek_{t+1})
\\ (1-\Discount\DeprFac) \ek_{t} & =  \ifTax{\TaxFree}{}\fFunc^{\kap}(\kap_{t}) - \ifTax{\kPriceAfterITC}{}\adj_{t}^{\kap}+\Delta \ek_{t+1}
\\ \ek_{t} & =  \frac{\ifTax{\TaxFree}{}\fFunc^{\kap}(\kap_{t}) - \ifTax{\kPriceAfterITC}{}\adj_{t}^{\kap}+\Delta \ek_{t+1}}{(1-\Discount\DeprFac)}
\end{aligned}\end{gathered}\end{equation}
and the phase diagram is constructed using the $\Delta \ek_{t+1} = 0$ locus.  In the vicinity of the steady state, we can assume $\adj_{t}^{\kap} \approx 0$ in which case the $\Delta \ek_{t+1}=0$ locus becomes
\begin{equation}\begin{gathered}\begin{aligned}
 \ek_{t} & =  \frac{\ifTax{\TaxFree}{}\fFunc^{\kap}(\kap_{t})}{(1-\Discount\DeprFac)} \label{eq:lamNearSS}
\end{aligned}\end{gathered}\end{equation}
which implies (since $\fFunc^{\kap}(\kap_{t})$ is downward sloping in $\kap_{t}$) that the $\Delta \ek_{t}=0$ locus
(that is, the $\lambda_{t}(\kap_{t})$ function that corresponds to $\Delta \ek_{t}=0$) is downward sloping.

The phase diagram is depicted below.% figure~\ref{fig:lPhaseDiag}.

\import {/Volumes/Data/Courses/Choice/LectureNotes/Investment/Handouts/EntrepreneurPF/LaTeX/}{lPhaseDiag.tex}

The steady state of the model will be the point at which $\kap_{t+1}=\kap_{t}=\check{\kap}$,
implying from \eqref{eq:kAccum} a steady-state investment rate of
\begin{equation}\begin{gathered}\begin{aligned}
  \check{\kap} & =  (\check{\kap}+\check{i})\DeprFac
\\ \check{i} & =  (1-\DeprFac)\check{\kap}/\DeprFac = (\depr/\DeprFac) \check{\kap}
\label{eq:iSS}
\end{aligned}\end{gathered}\end{equation}
and solving \eqref{kSSimplicit} for $ \ifTax{\TaxFree}{}\fFunc^{\kap}(\check{k})$
\begin{equation}\begin{gathered}\begin{aligned}
 \left(\frac{\ifTax{\kPriceAfterITC}{}(1 - \Discount\DeprFac)}{\Discount \DeprFac}\right) & =   \ifTax{\TaxFree}{}\fFunc^{\kap}(\check{k})
\end{aligned}\end{gathered}\end{equation}
which can be substituted into \eqref{eq:lamNearSS} to obtain the steady-state
value of $\ek$:
\begin{equation}\begin{gathered}\begin{aligned}
\check{\ek} & =   \left(\frac{\Rfree \ifTax{\kPriceAfterITC}{}}{\DeprFac}\right) .
\end{aligned}\end{gathered}\end{equation}

}{\AnswerOrBlankSpace{\clearpage\vfill\eject}{\vspace{0.0in}}}{}


We now wish to modify the problem in two ways.  First, we have been
assuming that the firm has only physical capital, and no financial
assets.  Second, we have been assuming that the manager running the
firm only cares about the PDV of profits; suppose instead we want to
assume that the firm is a small business run by an entrepreneur who
must live off the dividends of the firm, and thus they are maximizing
the discounted sum of utility from dividends $\utilFunc(\cRat_t)$
rather than just the level of discounted profits.  (Note that we
designate dividends by $\cRat_{t}$; dividends were not explicitly
chosen in the $\q$-model version of the problem, because the
Modigliani-Miller theorem says that the firm's value is unaffected by
its dividend policy).

We call the maximizer running this firm the `entrepreneur.'  The entrepreneur's level of monetary assets $m_{t}$ evolves according to
\begin{equation}\begin{gathered}\begin{aligned}
  m_{t+1} & =  \ifTax{\rev_{t+1}}{f_{t+1}}+\left(m_{t}-\ifTax{\xpend_{t}}{\inv_{t}-\adj_{t}}-\cRat_{t}\right)\Rfree.
\end{aligned}\end{gathered}\end{equation}

That is, next period the firm's money is next period's profits plus the
return factor on the money at the beginning of this period, minus
this period's investment and associated adjustment costs, minus dividends paid out
(which, having been paid out, are no longer part of the firm's money).


\ifthenelse{\boolean{Exam}}{
\item Show that t}{T}{he entrepreneur's Bellman equation can now be written}
\begin{equation}\begin{gathered}\begin{aligned}
  \vFunc_{t}(\kap_{t},m_{t}) & =  \max_{\{\inv_{t},\cRat_{t}\}}~~\utilFunc(\cRat_{t}) +\Discount \vFunc_{t+1}(\kap_{t+1},m_{t+1})\nonumber
\end{aligned}\end{gathered}\end{equation}

\answer{
Value is simply the discounted sum of utility from future dividends:
\begin{equation*}\begin{gathered}\begin{aligned}
  \vFunc_{t}(\kap_{t},m_{t}) & =  \max_{\{i,c\}_{t}^{\infty}} \sum_{n=0}^{\infty} \Discount^{n} \utilFunc(\cRat_{t+n})
\\ & =  \max_{\{i,c\}_{t}^{\infty}} \left(\utilFunc(\cRat_{t})+\Discount\sum_{n=0}^{\infty} \Discount^{n} \utilFunc(\cRat_{t+1+n})\right)
\\ & =  \max_{\{\inv_{t},\cRat_{t}\}}~~\utilFunc(\cRat_{t})+\Discount \vFunc_{t+1}(\kap_{t+1},m_{t+1}).
\end{aligned}\end{gathered}\end{equation*}

}{\AnswerOrBlankSpace{\clearpage\vfill\eject}{\vspace{0.0in}}}{}

Assume that $\fFunc$ and $\jFunc$ do not depend directly on $m_{t}$.  That is,
their partial derivatives with respect to $m_{t}$ are zero.



\ifthenelse{\boolean{Exam}}{
\item Use the first order condition with respect to dividends
and the Envelope theorem with respect to money to show that the
Euler equation for dividends is
\begin{equation}\begin{gathered}\begin{aligned}
 \uP(\cRat_{t}) & =  \Rfree\Discount \uP(\cRat_{t+1}) .
\end{aligned}\end{gathered}\end{equation}
}{  Then we will have}

\answer{
FOC wrt $\cRat_{t}$:
% Version where it's $\uP(d)$:
\begin{equation}\begin{gathered}\begin{aligned}
\uP(\cRat_{t}) & =  \Rfree\Discount \vNum^{m}_{t+1}
\end{aligned}\end{gathered}\end{equation}

Envelope wrt $m_{t}$:
\begin{equation}\begin{gathered}\begin{aligned}
  \vNum_{t}^{m} & =  \Rfree\Discount \vNum_{t+1}^{m}
\end{aligned}\end{gathered}\end{equation}
and combining the FOC with the Envelope theorem we get the usual
\begin{equation*}\begin{gathered}\begin{aligned}
  \vNum^{m}_{t} & =  \Rfree \Discount \vNum^{m}_{t+1}
\\ & =  \uP(\cRat_{t})
\\ & =  \Rfree\Discount \uP(\cRat_{t+1})
\\ & =  \uP(\cRat_{t+1})
\end{aligned}\end{gathered}\end{equation*}
where the last line follows because we have assumed $\Rfree\Discount=1$.
}{\AnswerOrBlankSpace{\clearpage\vfill\eject}{\vspace{0.0in}}}{}


\ifthenelse{\boolean{Exam}}{
\item Next explain why}{Now note that}
the value function can be rewritten as
\begin{equation}\begin{gathered}\begin{aligned}
  \vFunc_{t}(\kap_{t},m_{t}) & =  \max_{\{\inv_{t},m_{t+1}\}}~~\utilFunc((\ifTax{\rev_{t+1}}{f_{t+1}}-m_{t+1})/\Rfree+m_{t}-\ifTax{\xpend_t}{\inv_{t}-\adj_{t}}) +\Discount \vFunc_{t+1}(\kap_{t+1},m_{t+1}) \nonumber
\\ \label{eq:vOfmtp1}
\end{aligned}\end{gathered}\end{equation}


\answer{

This holds because maximizing with respect to $m_{t+1}$ (subject to the
accumulation equation) is equivalent to maximizing with respect to
the components of $m_{t+1}$.

}{\AnswerOrBlankSpace{\clearpage\vfill\eject}{\vspace{0.0in}}}{}

\ifthenelse{\boolean{Exam}}{
\item Now show that f}{F}or the version in \eqref{eq:vOfmtp1} the FOC with respect
to $\inv_{t}$ is
\begin{equation}\begin{gathered}\begin{aligned}
  \uP(\cRat_{t})(\ifTax{\kPriceAfterITC}{}(1+\adj_{t}^{\inv})-\ifTax{\TaxFree}{}f_{t+1}^{\kap}\DeprFac/\Rfree)& =  \DeprFac\Discount \vNum_{t+1}^{\kap} \label{eq:iFOCGen}
\end{aligned}\end{gathered}\end{equation}
\answer{

This holds because the derivative of the RHS of \eqref{eq:vOfmtp1}
with respect to $\inv_{t}$ is
\begin{equation}\begin{gathered}\begin{aligned}
  \uP(\cRat_{t})\left(\left(\frac{\partial \ifTax{\TaxFree}{}f_{t+1}}{\partial \kap_{t+1}}\frac{\partial \kap_{t+1}}{\partial \inv_{t}}\right)/\Rfree-\frac{\partial \ifTax{\kPriceAfterITC}{}\inv_{t}}{\partial \inv_{t}}-\frac{\partial \ifTax{\kPriceAfterITC}{}\adj_{t}}{\partial \inv_{t}}\right)+\Discount\left(\frac{\partial \kap_{t+1}}{\partial \inv_{t}}\right)\vFunc^{\kap}_{t+1}(\kap_{t+1},m_{t+1})
\end{aligned}\end{gathered}\end{equation}
(remember that $m_{t+1}$ is a control variable and thus its derivative with respect to investment is zero) so
the FOC translates to
\begin{equation}\begin{gathered}\begin{aligned}
  \uP(\cRat_{t})(\ifTax{\TaxFree}{}f_{t+1}^{\kap}\DeprFac/\Rfree-1-\ifTax{\kPriceAfterITC}{}\adj_{t}^{\inv})+\Discount\DeprFac \vNum_{t+1}^{\kap}&=0
\end{aligned}\end{gathered}\end{equation}
which reduces to \eqref{eq:iFOCGen}.
}{\null}{}



\ifthenelse{\boolean{Exam}}{\item Now
}{Now we can }use the envelope theorem with respect to $\kap_{t}$ to show that
\begin{equation}\begin{gathered}\begin{aligned}
  \vNum_{t}^{\kap} & =  \uP(\cRat_t)(\ifTax{\TaxFree}{}f_{t+1}^{\kap}\DeprFac/\Rfree-\ifTax{\kPriceAfterITC}{}\adj_{t}^{\kap})+\Discount \DeprFac \vNum_{t+1}^{\kap} \label{eq:kEnvelopeGen}
\end{aligned}\end{gathered}\end{equation}
\answer{

This can be seen by directly taking the derivative of the RHS of \eqref{eq:vOfmtp1} with respect to $\kap_{t}$:
\begin{equation}\begin{gathered}\begin{aligned}
  \uP(\cRat_{t})\left(\left(\frac{\partial \ifTax{\TaxFree}{}f_{t+1}}{\partial \kap_{t+1}}\frac{\partial \kap_{t+1}}{\partial \kap_{t}}\right)/\Rfree-\frac{\partial \ifTax{\kPriceAfterITC}{}\adj_{t}}{\partial \kap_{t}}\right)+ \Discount \left(\frac{\partial \kap_{t+1}}{\partial \kap_{t}}\right) \vNum_{t+1}^{\kap}
\end{aligned}\end{gathered}\end{equation}
and noting that the Envelope theorem tells us the derivatives with
respect to the controls $m_{t+1}$ and $\inv_{t}$ are zero while $\partial \kap_{t+1}/\partial \kap_{t} = \DeprFac$.

}{\AnswerOrBlankSpace{\clearpage\vfill\eject}{\vspace{0.0in}}}{}


\ifthenelse{\boolean{Exam}}{\item Next show how to
}{Now we can }combine \eqref{eq:iFOCGen} and \eqref{eq:kEnvelopeGen} to derive the Euler equation for investment
\begin{equation}\begin{gathered}\begin{aligned}
  \ifTax{\kPriceAfterITC}{} (1+\adj_{t}^{\inv}) & = \DeprFac\Discount \left[ \ifTax{\TaxFree}{}\fFunc^{\kap}(\kap_{t+1})+\ifTax{\kPriceAfterITC}{}(1+\adj_{t+1}^{\inv}-\adj_{t+1}^{\kap})\right] \label{eq:iEulerGen}
.
\end{aligned}\end{gathered}\end{equation}

\answer{

To see this, start with the Envelope theorem,
\begin{equation}\begin{gathered}\begin{aligned}
  \vNum_{t}^{\kap} & =  \uP(\cRat_{t})(\ifTax{\TaxFree}{}f_{t+1}^{\kap}\DeprFac/\Rfree-\ifTax{\kPriceAfterITC}{}\adj_{t}^{\kap})+\overbrace{\DeprFac \Discount \vNum_{t+1}^{\kap}}^{= \uP(\cRat_{t})(\ifTax{\kPriceAfterITC}{}(1+\adj_{t}^{\inv})-\ifTax{\TaxFree}{}f_{t+1}^{\kap}\DeprFac /\Rfree)~\mbox{from \eqref{eq:iFOCGen}}}
\\ & =  \uP(\cRat_{t})(\ifTax{\TaxFree}{}f_{t+1}^{\kap}\DeprFac/\Rfree-\ifTax{\kPriceAfterITC}{}\adj_{t}^{\kap})+\uP(\cRat_{t})(\ifTax{\kPriceAfterITC}{}(1+\adj_{t}^{\inv})-\ifTax{\TaxFree}{}f_{t+1}^{\kap}\DeprFac /\Rfree)
\\ & =  \uP(\cRat_{t})\ifTax{\kPriceAfterITC}{}\left(1+\adj_{t}^{\inv}-\adj_{t}^{\kap}\right) \label{eq:vkEqup}
\end{aligned}\end{gathered}\end{equation}
which means that we can rewrite \eqref{eq:iFOCGen} substituting the rolled-forward
version of \eqref{eq:vkEqup}
\begin{equation*}\begin{gathered}\begin{aligned}
  \uP(\cRat_{t})(\ifTax{\kPriceAfterITC}{}(1+\adj_{t}^{\inv})-\ifTax{\TaxFree}{}f_{t+1}^{\kap}\DeprFac/\Rfree)& =  \DeprFac\Discount \vNum_{t+1}^{\kap}
\\  & =  \DeprFac\Discount \uP(\cRat_{t+1})\ifTax{\kPriceAfterITC}{}\left(1+\adj_{t+1}^{\inv}-\adj_{t+1}^{\kap}\right)
\\   \ifTax{\kPriceAfterITC}{} (1+\adj_{t}^{\inv}) & = \DeprFac\Discount \left[ \ifTax{\TaxFree}{}\fFunc^{\kap}(\kap_{t+1})+\ifTax{\kPriceAfterITC}{}(1+\adj_{t+1}^{\inv}-\adj_{t+1}^{\kap})\right]
\end{aligned}\end{gathered}\end{equation*}
where the last line follows because with $\Rfree\Discount=1$ we know that $\cRat_{t+1}=\cRat_{t}$ implying
$\uP(\cRat_{t+1})=\uP(\cRat_{t})$.
}{\AnswerOrBlankSpace{\clearpage\vfill\eject}{\vspace{0.0in}}}{}


\ifthenelse{\boolean{Exam}}{\item Comment on the fact that the
  Euler equation for investment for the firm being run by a
  utility-maximizing manager, \eqref{eq:iEulerGen}, is identical to the
  Euler equation for the profit maximizing manager, \eqref{eq:iEuler},
  to discuss whether it matters, in this model, whether managers
  maximize profits or utility.  Similarly comment on whether there
  would be any evidence from consumption dynamics that the consumer
  was running a business with costly capital adjustment.
}{}


\answer{
  Since behavior (for either a firm manager or a consumer) is
  determined by Euler equations, and the Euler equations for both
  consumption and investment are identical in this model to the Euler
  equations for the standard models, there is no observable
  consequence for investment of the fact that the firm is being run by
  a utility maximizer, and there is no observable consequence for
  consumption of the fact that the consumer owns a business enterprise
  with costly capital adjustment.


}{\AnswerOrBlankSpace{\clearpage\vfill\eject}{\vspace{0.0in}}}{}



\ifthenelse{\boolean{Exam}}{
  \item

}{}

Now consider a firm of this kind that happens to have arrived in
period $t$ with positive monetary assets $m_{t}>0$ and with capital
equal to the steady-state target value $\kap_{t}=\check{k}$.

Suppose that a thief steals all the firm's monetary assets.

\ifthenelse{\boolean{Exam}}{Use the investment and consumption Euler equations
to show the consequences for monetary assets, capital, dividends, and investment subsequently.
}{
}

\answer{

The consequences for the firm are depicted below in figure~\ref{fig:mlossIRF}.

\import {/Volumes/Data/Courses/Choice/LectureNotes/Investment/Handouts/EntrepreneurPF/LaTeX/}{mlossIRF.tex}


Dividends follow a random walk.  Thus, there is a one-time downward adjustment
to the level of dividends to reflect the stolen money.  Thereafter dividends
are constant, as are monetary assets (which are constant at zero forever).

The theft of the money has no effect on investment or the capital stock,
because the firm's investment decisions are made on the basis of whether
they are profitable and the theft of the money has no effect on the profitability
of investments.

}{\AnswerOrBlankSpace{\clearpage\vfill\eject}{\vspace{0.0in}}}{}




\ifthenelse{\boolean{Exam}}{\item
}{}

Now consider another kind of shock: The firm's main building gets hit
by a meteor, destroying some of the firm's capital stock.
\ifthenelse{\boolean{Exam}}{Again show dynamics of monetary assets, capital, consumption,
and investment.}{}

\answer{
The results are depicted below in~\ref{fig:klossIRF}.

\import {/Volumes/Data/Courses/Choice/LectureNotes/Investment/Handouts/EntrepreneurPF/LaTeX/}{klossIRF.tex}

Again, because dividends follow a random walk, what the firm's managers
do is to assess the effect of the meteor shock on the firm's total value
and they adjust the level of dividends downward immediately to the sustainable
new level of dividends.  Thereafter there is no change in the level of
dividends.

Investment is more complicated.  The firm's capital stock is obviously reduced
below its steady-state value by the meteor, so there must be a period of
high investment expenditures to bring capital back toward its steady state.
However, the firm started out with monetary assets of zero.  Therefore the
high initial investment expenditures will be paid for by borrowing, driving
the firm's monetary assets to a permanent negative value (the firm goes into
debt to pay for its rebuilding).  Gradually over time the capital stock
is rebuilt back to its target level, and investment expenditures return
to zero (or the level consistent with replacing depreciated capital).

}{\clearpage}{\clearpage}



\ifthenelse{\boolean{Exam}}{\end{enumerate}}{}
\ifthenelse{\boolean{NumericalSolution}}{\input ./NumericalSolution.tex}{}% end if NumericalMethod

